\section{Related Work}
\label{sec:related-work}

\subsection{GPU-Accelerated Kalman Filtering}
\label{subsec:related-kalman}

GPU implementations of Kalman filtering expose parallelism at several different levels. Early work mapped the matrix operations within a filter step to CUDA and reported speedups when the matrices were large enough to amortize data movement and kernel-launch costs~\cite{huang2011gpu-kalman}. Application-specific systems more often exploit independent estimation problems. The GPU track-fitting implementation of Ai et al.~\cite{ai2021gpu-kalman} processes many particle tracks concurrently, while FastTrack parallelizes the prediction and update of multiple tracked objects and integrates the filter with the surrounding tracking pipeline~\cite{liu2024fasttrack}. These designs are effective when a workload supplies many simultaneous filters, although they do not remove the recurrence or fine-grained execution overhead of one long, small-state trajectory.

Another line of work derives temporally parallel filtering and smoothing algorithms. S\"arkk\"a and Garc\'ia-Fern\'andez formulate Bayesian smoothing operations as associative elements, enabling prefix-scan algorithms with logarithmic span~\cite{sarkka2021temporal}. Such reformulations are attractive when the model admits associative composition and minimizing single-trajectory latency is the primary goal. The downhole traction robot positioning simulation studied here retains the original nonlinear SINS propagation, measurement selection, error feedback, and covariance recursion in their program order. We shorten this sequential path by fusing the complete scan and keeping its recurrent state local, then use independent simulation trajectories---rather than time steps from one trajectory---to occupy the GPU.

\subsection{MegaKernel and Kernel Fusion}
\label{subsec:related-fusion}

Kernel fusion combines operations so that intermediate values can remain on chip and repeated kernel launches can be avoided. MegaKernels enlarge this scope to a substantial computation region; for example, Jia et al.~\cite{jia2020megakernels} place the stages of fast convolution in one GPU kernel to reuse intermediate data and reduce launch and synchronization overhead. Recent compiler and systems research has extended fusion beyond manually selected, memory-bound operator chains. Welder constructs a tile graph to jointly search fusion and tile schedules according to memory traffic~\cite{shi2023welder}. Chimera uses an analytical model to select fusion plans for compute-intensive operators~\cite{zheng2023chimera}, and MCFuser combines memory-bound operators with compute-intensive operators while controlling recomputation and synchronization costs~\cite{zhang2024mcfuser}. These systems show that profitable fusion depends on both the operator graph and the hardware cost created by a candidate fused region.

More recent systems search larger transformation spaces. Mirage represents tensor programs at kernel, thread-block, and thread levels, allowing its superoptimizer to combine algebraic transformations, schedules, and custom-kernel generation~\cite{wu2025mirage}. STOF uses compilation templates and a two-stage search to select adaptive fusion strategies for sparse Transformer inference~\cite{dai2026stof}. These approaches primarily optimize tensor graphs whose operators retain substantial spatial parallelism. Our target presents a different fusion boundary: a complete sequential SINS/EKF trajectory is compiled as one Triton program~\cite{triton}, and its small navigation state and covariance remain program-local across time steps. Parallelism is recovered across independent trajectories, while fusion reduces the latency of each recurrent scan. This specialization complements graph-level fusion systems by addressing a long-lived kernel with an unavoidable loop-carried dependency.

\subsection{Mixed Precision for Scientific Computing}
\label{subsec:related-precision}

Mixed-precision methods have been widely studied in scientific computing and GPU acceleration~\cite{abdelfattah2021survey}. Baboulin et al.~\cite{mixed-precision} combine lower-precision kernels with higher-precision residual evaluation and correction in numerical linear algebra. Haidar et al.~\cite{haidar2018tensorcores} further map low-precision factorization and solves to GPU Tensor Cores while retaining higher precision for accuracy-critical operations.

Automatic precision selection provides a complementary approach. Precimonious searches variable- and operation-level assignments under an application-provided error constraint~\cite{rubio2013precimonious}. CoMP performs operator-wise precision selection for neural-network training through convergence-aware and performance-oriented sampling~\cite{dai2025comp}. Its search also considers format-conversion overhead and the availability of Tensor Core execution, illustrating the importance of evaluating precision plans against both numerical requirements and the characteristics of the target GPU. Unlike these per-operation or per-variable search schemes, \textsc{Trident} fixes one compile-time precision plan per device and validates it over a complete recursive trajectory, so rounding that only emerges over a long horizon is measured before the plan is adopted.
